\documentclass[11pt,a4paper]{article}
% ============================================================
%  Structural Persistence Series — Shared Preamble
% ============================================================
\usepackage{fontspec}
\setmainfont{Hiragino Mincho ProN}
\setsansfont{Hiragino Sans}
\setmonofont{Menlo}

\usepackage{iftex}
\ifXeTeX
  \XeTeXlinebreaklocale "ja"
  \XeTeXlinebreakskip=0pt plus 1pt minus 0.1pt
\fi

\usepackage[margin=22mm]{geometry}
\usepackage{amsmath,amssymb,amsthm}
\usepackage{xcolor}
\usepackage{graphicx}
\usepackage{hyperref}
\usepackage{booktabs}
\usepackage{fancyhdr}
% enumitem not available in this TeX installation

% --- Colours ------------------------------------------------
\definecolor{PaperAccent}{HTML}{1F4E79}
\definecolor{PaperGray}{HTML}{51606F}

\hypersetup{
  colorlinks=true,
  linkcolor=PaperAccent,
  urlcolor=PaperAccent,
  citecolor=PaperAccent,
  pdflang=ja
}
\urlstyle{same}

% --- Typography ---------------------------------------------
\setlength{\parindent}{1.5em}
\setlength{\parskip}{0pt}
\setlength{\emergencystretch}{3em}
\linespread{1.08}
\setlength{\abovedisplayskip}{8pt plus 2pt minus 4pt}
\setlength{\belowdisplayskip}{8pt plus 2pt minus 4pt}
\setlength{\abovedisplayshortskip}{6pt plus 2pt minus 2pt}
\setlength{\belowdisplayshortskip}{6pt plus 2pt minus 2pt}

% --- Section label names ------------------------------------
\renewcommand{\abstractname}{要旨}

% --- Theorem-like environments ------------------------------
\theoremstyle{plain}
\newtheorem{proposition}{命題}[section]
\newtheorem{lemma}[proposition]{補題}
\newtheorem{corollary}[proposition]{系}

\theoremstyle{definition}
\newtheorem{definition}{定義}[section]
\newtheorem{assumption}{条件}[section]

\theoremstyle{remark}
\newtheorem*{remark}{注}

% --- Running header -----------------------------------------
\pagestyle{fancy}
\fancyhf{}
\renewcommand{\headrulewidth}{0.4pt}
\fancyhead[L]{\small\textcolor{PaperGray}{\nouppercase{\leftmark}}}
\fancyhead[R]{\small\textcolor{PaperGray}{\thepage}}
\fancypagestyle{plain}{%
  \fancyhf{}
  \renewcommand{\headrulewidth}{0pt}
  \fancyfoot[C]{\small\textcolor{PaperGray}{\thepage}}
}

% --- Title block macro --------------------------------------
\newcommand{\PaperTitleBlock}[3]{%
  % #1 = title, #2 = subtitle, #3 = date
  \thispagestyle{plain}%
  \begin{center}
  {\LARGE\bfseries #1\par}
  \vspace{0.4em}
  {\normalsize #2\par}
  \vspace{1.0em}
  {\large Akihito Sunagawa\par}
  \vspace{0.3em}
  {\normalsize #3\par}
  \end{center}
  \vspace{1.6em}
}

% Legacy 2-arg version (backward compat)
\newcommand{\PaperTitleBlockSimple}[2]{%
  \PaperTitleBlock{#1}{}{#2}
}


\hypersetup{pdflang=en}
\renewcommand{\abstractname}{Abstract}

\begin{document}
\linespread{1.12}\selectfont
\thispagestyle{plain}
\begin{center}
{\Large\bfseries Structural Persistence Theory\par}
\vspace{0.3em}
{\normalsize English Abstract\par}
\vspace{0.5em}
{\small Akihito Sunagawa \quad April 16, 2026\par}
\end{center}
\vspace{0.8em}

\begin{abstract}
This note is the shortest public English entry point to the project.

The central hypothesis is that two problems usually discussed separately may share the same structural mechanism: reasoning degradation in long conversations, and catastrophic forgetting under continual learning. Both can be read as cases where the set of states that can still preserve a target structure gradually shrinks as unresolved contradictions or premise-changing updates accumulate.

At the theoretical core, structural loss is defined by the log-ratio of successive shrinkage in the set of states that can still sustain the structure. If a system starts with a structure-preserving state set
\[
V^{(0)} \supseteq V^{(1)} \supseteq \cdots \supseteq V^{(n)},
\]
and the stage loss is defined as
\[
l_i = -\ln \frac{m(V^{(i)})}{m(V^{(i-1)})},
\]
then the remaining survivable region takes the exponential form
\[
m(V^{(n)}) = m(V^{(0)}) e^{-L}, \qquad L = \sum_i l_i.
\]
In this framework, the exponential form is not an extra empirical assumption; it follows directly from representing cumulative structural loss as successive multiplicative shrinkage measured in log-ratio form.

Empirically, the inference experiments suggest that long-context degradation is not mainly a length problem, but a contradiction-management problem: externally organizing contradictory updates preserves coherence better than leaving the same contradictions unresolved. The continual-learning experiments suggest that LoRA-style sequential updating behaves more like overwrite than clean accumulation when dependency-linked knowledge must be reorganized after premise changes. Dependency-aware replay and adapter separation help, but do not fully solve high-fidelity long-term retention.

The broader architectural implication is that long-horizon intelligence may require more than prompt engineering or more training alone. The work points toward external contradiction metabolism, multi-layer memory, premise-dependent reorganization, rollbackable state management, and persistent internal-model maintenance across time.
\end{abstract}

\end{document}
